\fichatitulo{50 · RECUPERAÇÃO JURÍDICA}{Workshop NLLP · 2025}{Document-Level Retrieval Mismatch}
{\small Reuter et al.. \textit{Document-Level Retrieval Mismatch}. Workshop NLLP · 2025. \href{https://aclanthology.org/2025.nllp-1.3/}{Fonte consultada}.\par}

\campo{Problema e objetivo}
Que tipo de falha ocorre quando o recuperador seleciona o documento errado em bases legais semelhantes?

\campo{Principal contribuição · síntese interpretativa}
Define e quantifica DRM e testa Summary-Augmented Chunking para preservar contexto global no chunk.

\campo{Método / natureza da evidência}
Experimentos em tarefas jurídicas de recuperação, comparando baselines e a técnica SAC.

\campo{Resultado ou argumento principal}
SAC reduz DRM e melhora precisão e recall no nível textual nas condições estudadas; evidencia que recuperação correta do trecho depende da identificação do documento.

\campo{Excertos-chave · seleção literal}
\begin{itemize}
\excerto{1}{the retriever selects information from entirely incorrect source documents}{Resumo.}
\end{itemize}

\campo{Limitações e análise crítica}
A técnica foi testada em dados jurídicos; não se deve presumir ganho no Senado sem ablação e validação no corpus parlamentar. Resumos sintéticos também introduzem uma transformação que precisa ser documentada.

\campo{Aproveitamento na dissertação · proposta}
Metodologia: medir erro de documento separadamente do erro de trecho e registrar efeitos do chunking.

\registro{Fonte consultada nesta preparação: PDF publicado; consulta focal do resumo, método, resultados e conclusão. Chave: \texttt{reuterEtAl2025reliableRetrieval}.}
